% =====================================================================
% Appendix: constructing and scoring the CATE / ATE predictive density
% Self-contained; \input into the ICLR main file, or compile standalone
% via main_standalone.tex in this folder.
% =====================================================================
\section{Constructing the treatment-effect predictive distribution}
\label{app:density}

Every method we compare emits a \emph{discrete} predictive distribution over
outcomes rather than a point estimate, so the interval diagnostics in
Section~\ref{sec:calibration} are computed from the models' own predictive
objects with no post-hoc smoothing, resampling, or density estimation. This
appendix states the construction exactly, because the methods differ in what
they emit --- one marginal per arm versus a joint over both --- and those two
cases require different reductions to a distribution over
$\tau = Y^{(1)} - Y^{(0)}$.

\subsection{Notation}
\label{app:notation}

Fix a query covariate vector $x$. Write $Y^{(t)}$ for the potential outcome
under $\mathrm{do}(T=t)$, $t \in \{0,1\}$, and $\tau = Y^{(1)} - Y^{(0)}$ for
the individual treatment effect. All methods discretise the outcome axis with
a \emph{bar distribution}: a set of $K$ bins defined by edges
\begin{equation}
  e_0 < e_1 < \cdots < e_K ,
  \qquad
  w_k = e_k - e_{k-1},
  \qquad k = 1,\dots,K ,
\end{equation}
together with a probability vector $p^{(t)} = (p^{(t)}_1,\dots,p^{(t)}_K)$,
$p^{(t)}_k \ge 0$, $\sum_k p^{(t)}_k = 1$, obtained as a softmax of the
model's logits. Each bin is summarised by a \emph{representative point}
$m_k$; the choice of $m_k$ is model-specific and is the subject of
Section~\ref{app:models}. For a uniform grid with $w_k \equiv \Delta$ the
natural choice is the midpoint $m_k = \tfrac12(e_{k-1} + e_k)$, so that
\begin{equation}
  m_k = m_1 + (k-1)\Delta .
  \label{eq:lattice}
\end{equation}
Equation~\eqref{eq:lattice} --- that the representative points lie on a
lattice of spacing $\Delta$ --- is what makes the reductions below exact; we
return to its failure in Section~\ref{app:nonuniform}.

We write $\mathcal{M} = \{m_1,\dots,m_K\}$ for the support and
$\hat F^{(t)}$ for the implied CDF.

\subsection{Two-dimensional heads: anti-diagonal projection}
\label{app:joint}

A 2D head emits a joint law over the pair of potential outcomes: a matrix
$P = (P_{ij})_{i,j=1}^{K}$ with $P_{ij} \ge 0$, $\sum_{ij} P_{ij} = 1$, where
\begin{equation}
  P_{ij} \;=\; \Pr\!\left[\, Y^{(0)} \in B_i ,\; Y^{(1)} \in B_j \,\right],
  \qquad B_k = [e_{k-1}, e_k) .
\end{equation}
Assigning each bin its representative point gives $\tau = m_j - m_i$, which
by \eqref{eq:lattice} equals $(j-i)\Delta$. The induced law of $\tau$ is
therefore supported on the $2K-1$ lattice points $\{d\Delta\}_{d=-(K-1)}^{K-1}$
with
\begin{equation}
  \boxed{\;
  \Pr[\tau = d\Delta] \;=\; \sum_{\substack{i,j \\ j - i = d}} P_{ij}
  \;}
  \label{eq:diag}
\end{equation}
i.e.\ the sum of the $d$-th anti-diagonal of $P$. No independence assumption
is made: \eqref{eq:diag} uses whatever dependence between the arms the model
represents, which is precisely what a 2D head is for.

\subsection{One-dimensional heads: coupling the two arms}
\label{app:marginals}

A 1D head emits only the two marginals $p^{(0)}, p^{(1)}$. These do not
determine the law of $\tau$: any coupling $\pi$ with the given marginals is
consistent with them, and different couplings give different answers. A
choice must therefore be made explicit.

\paragraph{Independence coupling.}
Taking $\pi_{ij} = p^{(0)}_i p^{(1)}_j$ and substituting into
\eqref{eq:diag},
\begin{equation}
  \boxed{\;
  \Pr[\tau = d\Delta]
  \;=\; \sum_{\substack{i,j \\ j-i=d}} p^{(0)}_i\, p^{(1)}_j
  \;=\; \bigl(p^{(1)} * \widetilde{p^{(0)}}\bigr)_d ,
  \qquad \widetilde{p^{(0)}}_k = p^{(0)}_{K+1-k} ,
  \;}
  \label{eq:conv}
\end{equation}
a discrete convolution of $p^{(1)}$ with the reflection of $p^{(0)}$. We
evaluate \eqref{eq:conv} directly for small $K$ and by FFT above a threshold
$K_{\mathrm{fft}}$, at which point the cost is $O(K \log K)$ rather than
$O(K^2)$.

\paragraph{Comonotonic coupling.}
The opposite extreme couples the arms through their quantile functions,
\begin{equation}
  \tau(u) \;=\; \bigl(\hat F^{(1)}\bigr)^{-1}(u) - \bigl(\hat F^{(0)}\bigr)^{-1}(u),
  \qquad u \sim \mathrm{Unif}(0,1),
\end{equation}
which is the coupling maximising the correlation between the arms and hence
minimising $\operatorname{Var}(\tau)$.

\paragraph{Why the choice matters.}
Let $\rho$ denote the correlation between $Y^{(0)}$ and $Y^{(1)}$ under the
data-generating process. Then
\begin{equation}
  \operatorname{Var}(\tau)
  = \sigma_0^2 + \sigma_1^2 - 2\rho\,\sigma_0\sigma_1 ,
  \label{eq:vartau}
\end{equation}
so the independence coupling ($\rho = 0$) overstates the spread of $\tau$
whenever $\rho > 0$, by a factor of
$\bigl[(\sigma_0^2+\sigma_1^2)/(\sigma_0^2+\sigma_1^2-2\rho\sigma_0\sigma_1)\bigr]^{1/2}$
in standard deviation. On generators that draw both potential outcomes from
shared exogenous noise, $\rho$ is close to $1$ and this factor is large. We
report the independence coupling as the default because it is the assumption
a 1D head implicitly licenses; the comparison to a 2D head on the same data
is then a measurement of what the joint representation buys, and
\eqref{eq:vartau} is the reason a 2D head can be sharper without being
over-confident.

\subsection{Outcome standardisation}
\label{app:standardisation}

Models predict on a standardised axis. Let arm $t$ use the affine map
\begin{equation}
  \widetilde{y} = \frac{y - \mu_t}{\sigma_t},
  \qquad\text{equivalently}\qquad
  y = \sigma_t \widetilde{y} + \mu_t .
\end{equation}
Then
\begin{equation}
  \tau
  = \bigl(\sigma_1 \widetilde{Y}^{(1)} + \mu_1\bigr)
  - \bigl(\sigma_0 \widetilde{Y}^{(0)} + \mu_0\bigr) .
  \label{eq:perarm}
\end{equation}
Two regimes follow.

\paragraph{Pooled ($\mu_0=\mu_1=\mu$, $\sigma_0=\sigma_1=\sigma$).}
Equation~\eqref{eq:perarm} collapses to
$\tau = \sigma\,(\widetilde{Y}^{(1)} - \widetilde{Y}^{(0)})$: the location
parameter cancels and a single scalar $\sigma$ maps the standardised effect
to raw units. Both reductions above then apply verbatim on the standardised
axis, with atoms multiplied by $\sigma$.

\paragraph{Per-arm ($\mu_0 \neq \mu_1$ or $\sigma_0 \neq \sigma_1$).}
The location parameters do \emph{not} cancel, and no single scalar recovers
$\tau$ from $\widetilde{Y}^{(1)} - \widetilde{Y}^{(0)}$. Applying the pooled
formula regardless omits $\mu_1 - \mu_0$ and mis-scales one arm --- a pure
translation of the $\tau$ density, which leaves its width intact while moving
its centre, and therefore degrades coverage without affecting the spread
diagnostics. When a method is evaluated under per-arm standardisation we map
both arms to raw units first (Section~\ref{app:nonuniform}) and only then
couple them.

\subsection{Non-uniform bins}
\label{app:nonuniform}

Both reductions rely on \eqref{eq:lattice}. When the edges are not equally
spaced --- as for quantile-spaced borders --- the differences $m_j - m_i$ do
not lie on a lattice and neither \eqref{eq:diag} nor \eqref{eq:conv} applies.

We therefore re-express such a predictive on a uniform grid before coupling.
Treating the mass of each bin as uniformly distributed over that bin (the
definition of a bar distribution) gives the piecewise-linear CDF
\begin{equation}
  \hat F^{(t)}(e_k) = \sum_{k' \le k} p^{(t)}_{k'} ,
\end{equation}
and for a target grid $e'_0 < \cdots < e'_{K'}$ the re-expressed masses are
\begin{equation}
  \boxed{\;
  q^{(t)}_{k} \;=\; \hat F^{(t)}(e'_{k}) - \hat F^{(t)}(e'_{k-1}) ,
  \qquad k = 1,\dots,K' .
  \;}
  \label{eq:rebin}
\end{equation}
This preserves total mass exactly and the mean to the resolution of the
target grid. Two properties of the target grid matter.

\begin{enumerate}
\item \textbf{Both arms must share it.} If each arm were given its own grid
      $\{e'^{(0)}\}$, $\{e'^{(1)}\}$ with different origins, the index
      difference $j-i$ used by \eqref{eq:diag}--\eqref{eq:conv} would no
      longer correspond to $m_j - m_i$; the offset between the two origins
      would be absorbed silently into $\tau$ as a translation. We take
      $e'_0 = \min_t \min_k e^{(t)}_k$ and
      $e'_{K'} = \max_t \max_k e^{(t)}_k$.
\item \textbf{It must cover the full support.} Truncating the grid and
      clamping out-of-range mass onto the endpoints moves that mass, biasing
      $\mathbb{E}[\tau]$. Resolution is instead obtained by increasing $K'$,
      chosen so the grid spacing is comparable to the median inter-bin gap of
      the source and capped at $K'_{\max}$ for tractability. $K'_{\max}$ is
      held fixed across all benchmarks, so a given method is always scored at
      one resolution.
\end{enumerate}

\subsection{Model-specific predictive objects}
\label{app:models}

\begin{table}[h]
\centering
\caption{What each method emits and how $\tau$ is obtained. $K$ denotes the
method's own bin count; no method is re-discretised to another method's grid.}
\label{tab:heads}
\small
\begin{tabular}{llll}
\toprule
Method & Emits & Grid & $\tau$ via \\
\midrule
CausalPFN-1D & two marginals & uniform & \eqref{eq:conv} \\
CausalPFN-2D & joint $P$ & uniform & \eqref{eq:diag} \\
Graph2D & joint $P$ & uniform & \eqref{eq:diag} \\
UWYK-1D & two marginals & uniform & \eqref{eq:conv} \\
DoPFN-bb & joint $P$ & uniform & \eqref{eq:diag} \\
DoPFN (native) & two marginals & non-uniform & \eqref{eq:rebin} then \eqref{eq:conv} \\
\bottomrule
\end{tabular}
\end{table}

\paragraph{Uniform-grid heads.}
CausalPFN, Graph2D and UWYK place their bins on an evenly spaced grid, so
$m_k$ is the bin midpoint and \eqref{eq:lattice} holds directly. UWYK's head
additionally carries two half-normal tail components beyond the outermost
edges; we treat them as ordinary boundary bins, which is the same
approximation the method's own resolution-matched variant makes.

\paragraph{Full-support bar distributions.}
DoPFN's native head uses a full-support bar distribution: the interior bins
are finite, while the first and last carry half-normal tails extending beyond
the grid. Writing $\mathrm{HN}(s)$ for a half-normal with scale $s$, the tail
scales are set so that a fixed probability $\eta$ falls within the outermost
bin width,
\begin{equation}
  s_{\mathrm{L}} = \frac{w_1}{F^{-1}_{\mathrm{HN}(1)}(\eta)},
  \qquad
  s_{\mathrm{R}} = \frac{w_K}{F^{-1}_{\mathrm{HN}(1)}(\eta)},
\end{equation}
and since $\mathbb{E}[\mathrm{HN}(s)] = s\sqrt{2/\pi}$ the representative
points of the two boundary bins are
\begin{equation}
  m_1 = e_1 - s_{\mathrm{L}}\sqrt{2/\pi},
  \qquad
  m_K = e_{K-1} + s_{\mathrm{R}}\sqrt{2/\pi},
  \label{eq:tailmeans}
\end{equation}
which lie \emph{outside} their bins. Using midpoints for these two entries
mis-states $\mathbb{E}[Y^{(t)}]$ whenever the tails carry non-negligible mass.

\paragraph{Consistency between the point estimate and the density.}
A calibration table is interpretable only if the reported point estimate and
the reported intervals describe the same predictive object. We therefore
define each method's $\tau$ density using the \emph{same} bin representatives
that its own point estimate uses. Concretely, if a released implementation
reports
\begin{equation}
  \hat\tau(x) \;=\; \sum_k p^{(1)}_k m^{\star}_k - \sum_k p^{(0)}_k m^{\star}_k
\end{equation}
for some representatives $m^{\star}$, then $m^{\star}$ --- not the midpoints
--- is what enters \eqref{eq:conv}. This matters when $m^{\star}$ differs from
the midpoint: a representative computed on a different scale from the edges,
say $m^{\star}_k = e_{k-1} + w_k/(2\alpha)$ for a normalisation constant
$\alpha \neq 1$, shifts each bin by $\tfrac{w_k}{2}(\alpha^{-1}-1)$. For a
uniform grid this is a constant that cancels in $\tau$; for a non-uniform
grid it does not, and the residual is
\begin{equation}
  \hat\tau_{\star} - \hat\tau_{\mathrm{mid}}
  = \frac{\alpha^{-1}-1}{2} \sum_k \bigl(p^{(1)}_k - p^{(0)}_k\bigr) w_k ,
  \label{eq:repdiff}
\end{equation}
which scales with the spread of the bin widths and with how far $\alpha$ is
from unity. We adopt each method's own $m^{\star}$ so that PEHE and the
interval scores are never computed from different estimators, and note
\eqref{eq:repdiff} as the magnitude of the difference this choice makes.

\subsection{From CATE to ATE}
\label{app:ate}

For a dataset with queries $x_1,\dots,x_n$ and per-query effect densities
$\nu_1,\dots,\nu_n$ on a common support, the predictive for the average
treatment effect is their $2$-Wasserstein barycenter,
\begin{equation}
  \bar\nu \;=\; \argmin_{\nu} \; \frac{1}{n}\sum_{q=1}^{n} W_2^2(\nu, \nu_q) .
\end{equation}
In one dimension this has the closed form
\begin{equation}
  \boxed{\;
  \bar F^{-1}(u) \;=\; \frac{1}{n}\sum_{q=1}^{n} F_q^{-1}(u),
  \qquad u \in (0,1),
  \;}
\end{equation}
the pointwise average of the quantile functions. The barycenter is scored
against the dataset's realised ATE, $\bar\tau = \tfrac{1}{n}\sum_q \tau^\star_q$.
Averaging quantile functions rather than densities preserves the shape of the
per-query predictive; averaging densities pointwise would instead produce a
mixture, whose spread reflects heterogeneity across queries rather than
uncertainty about the average.

\section{Interval diagnostics}
\label{app:metrics}

Let $\nu$ be a predictive for $\tau$ with CDF $F$, and let $\tau^\star$ be the
realised effect. For a nominal level $1-\alpha$ define
\begin{equation}
  \ell = F^{-1}(\alpha/2), \qquad u = F^{-1}(1-\alpha/2).
\end{equation}

\paragraph{Coverage and length.}
$\mathrm{Cov} = \mathbb{1}\{\ell \le \tau^\star \le u\}$ and
$\mathrm{Len} = u - \ell$, averaged over queries. Neither is interpretable
alone: a sufficiently wide interval attains nominal coverage at any accuracy.

\paragraph{Interval score.}
The two are priced jointly by the interval (Winkler) score
\begin{equation}
  \boxed{\;
  \mathrm{IS}_{\alpha}(\ell,u;\tau^\star)
  = (u - \ell)
  + \frac{2}{\alpha}(\ell - \tau^\star)\,\mathbb{1}\{\tau^\star < \ell\}
  + \frac{2}{\alpha}(\tau^\star - u)\,\mathbb{1}\{\tau^\star > u\} ,
  \;}
\end{equation}
which is proper, equals $\mathrm{Len}$ whenever the interval covers, and
otherwise adds $2/\alpha$ times the miss distance. We report
$\mathrm{IS}_{\alpha}$ at the same $\alpha$ as the coverage and length
columns, so all three describe one interval.

\paragraph{Relation to CRPS.}
Averaging $\mathrm{IS}_{\alpha}$ over a grid of levels
$\alpha_1 > \cdots > \alpha_J$ together with a median term gives the weighted
interval score, which approximates
\begin{equation}
  \mathrm{CRPS}(F, \tau^\star)
  = \int_{-\infty}^{\infty}\bigl(F(z) - \mathbb{1}\{z \ge \tau^\star\}\bigr)^2 \, dz .
\end{equation}
For a discrete $F$ with atoms $m_1 < \cdots < m_K$ and masses $p_k$ this
integral is evaluated in closed form as a sum over the intervals between
consecutive atoms. We report $\mathrm{IS}_{\alpha}$ rather than CRPS or WIS
as the headline number because it scores the interval actually being
reported; CRPS and WIS are retained in the released artefacts.

\paragraph{Spread diagnostic.}
We additionally report
$\mathrm{sd\text{-}ratio} = \operatorname{sd}(\nu) / \operatorname{sd}(\tau^\star)$,
the ratio of mean predictive standard deviation to the standard deviation of
the true effects. This is a units check rather than a quality measure: a
predictive expressed on the wrong scale still receives a finite score, and
$\mathrm{sd\text{-}ratio}$ far from unity identifies that failure mode, which
the other three columns cannot distinguish from genuine over- or
under-dispersion.
